% ===========================================================================
% PROYECTO FINAL · Compiladores, FI UNAM
% Diseño e implementación de un analizador léxico y sintáctico para MiniC
% ===========================================================================
\documentclass[12pt]{article}

\usepackage[utf8]{inputenc}
\usepackage[T1]{fontenc}
\usepackage{lmodern}
\usepackage[spanish,es-nodecimaldot,es-noshorthands,es-tabla]{babel}
\usepackage{microtype}
\usepackage{textcomp}

\usepackage{graphicx}
\usepackage{geometry}
\usepackage{setspace}
\usepackage{array}
\usepackage{tabularx}
\usepackage{booktabs}
\usepackage{longtable}
\usepackage[table]{xcolor}
\usepackage{enumitem}
\usepackage[most]{tcolorbox}
\usepackage{titlesec}
\usepackage{fancyhdr}
\usepackage{float}
\usepackage{pdflscape}
\usepackage{needspace}
\usepackage{amsmath,amssymb}
\usepackage{listings}
\usepackage{tikz}
\usetikzlibrary{automata,positioning,arrows.meta,shapes.geometric}
\usepackage[titles]{tocloft}
\setlength{\cftsecnumwidth}{2em}
\setlength{\cftsubsecnumwidth}{2.8em}
\usepackage{caption}
\captionsetup{justification=centering}
\usepackage{xurl}
\usepackage[hidelinks]{hyperref}

% -- Datos del curso y del trabajo -------------------------------------------
\newcommand{\Materia}{Compiladores}
\newcommand{\Semestre}{2027-1}
\newcommand{\Profesor}{M.I. Manuel Enrique Castañeda Castañeda}
\newcommand{\Grupo}{02}
\newcommand{\FechaEntrega}{14 de noviembre de 2026}
\newcommand{\TituloTrabajo}{Proyecto final}
\newcommand{\SubtituloTrabajo}{Analizador léxico y sintáctico para MiniC}
% ---------------------------------------------------------------------------

\definecolor{grisclaro}{gray}{0.955}
\definecolor{grismedio}{gray}{0.88}
\definecolor{grisborde}{gray}{0.70}

\geometry{a4paper, top=3.0cm, bottom=2.6cm, left=2.8cm, right=2.8cm,
headheight=15pt, headsep=0.7cm, footskip=1.1cm}
\setstretch{1.15}
\setlength{\parskip}{0.45em}
\setlength{\parindent}{0pt}
\renewcommand{\topfraction}{0.9}
\renewcommand{\bottomfraction}{0.8}
\renewcommand{\textfraction}{0.07}
\renewcommand{\floatpagefraction}{0.75}

\pagestyle{fancy}
\fancyhf{}
\renewcommand{\headrulewidth}{0.4pt}
\fancyhead[L]{\footnotesize\textbf{\Materia}\ \textbar\ Grupo \Grupo}
\fancyhead[R]{\footnotesize\TituloTrabajo}
\fancyfoot[C]{\footnotesize\color{black!60}\thepage}
\fancypagestyle{plain}{\fancyhf{}\renewcommand{\headrulewidth}{0pt}%
  \fancyfoot[C]{\footnotesize\color{black!60}\thepage}}

\titleformat{\section}
{\Large\bfseries}{\thesection.}{0.6em}{}[{\titlerule[1.0pt]}]
\titlespacing*{\section}{0pt}{1.6em}{0.9em}

\titleformat{\subsection}
{\large\bfseries}{\thesubsection.}{0.6em}{}
\titlespacing*{\subsection}{0pt}{1.2em}{0.5em}

\newcommand{\cmd}[1]{\texttt{\small #1}}
\newcommand{\nt}[1]{\textit{#1}}
\newcommand{\eps}{$\varepsilon$}

\newtcolorbox{apartado}[1]{
enhanced,
colback=white, colframe=black,
boxrule=0.6pt, arc=0pt,
left=9pt, right=9pt, top=7pt, bottom=7pt,
fonttitle=\bfseries\small, coltitle=white, colbacktitle=black,
attach boxed title to top left={xshift=8pt, yshift=-\tcboxedtitleheight/2},
boxed title style={arc=1pt, boxrule=0pt},
title={#1},
top=12pt
}

\newtcolorbox{respuesta}{
enhanced, breakable,
colback=grisclaro, colframe=black,
boxrule=0pt, leftrule=2.5pt, arc=0pt,
left=10pt, right=6pt, top=6pt, bottom=6pt
}

\lstdefinestyle{codigo}{
basicstyle=\ttfamily\footnotesize,
columns=fullflexible, keepspaces=true,
backgroundcolor=\color{grisclaro},
frame=leftline, framerule=2.5pt, rulecolor=\color{black},
xleftmargin=8pt, framexleftmargin=6pt, framesep=6pt,
aboveskip=0.8em, belowskip=0.6em,
keywordstyle=\bfseries, commentstyle=\color{black!60}\itshape,
showstringspaces=false, breaklines=true,
literate={á}{{\'a}}1 {é}{{\'e}}1 {í}{{\'i}}1 {ó}{{\'o}}1 {ú}{{\'u}}1 {ñ}{{\~n}}1
         {Á}{{\'A}}1 {É}{{\'E}}1 {Í}{{\'I}}1 {Ó}{{\'O}}1 {Ú}{{\'U}}1
         {→}{{$\rightarrow$}}1 {ε}{{$\varepsilon$}}1
}
\lstdefinestyle{fuente}{
style=codigo, basicstyle=\ttfamily\scriptsize, numbers=left, numberstyle=\tiny\color{black!55}, numbersep=12pt,
xleftmargin=22pt, framexleftmargin=20pt, language=C
}
\lstset{style=codigo}
\renewcommand{\lstlistingname}{Código}

% Los logos se toman de la carpeta IMAGENES (si no existen, se deja el espacio)
\newcommand{\logo}[1]{\IfFileExists{#1}{\includegraphics[height=2.9cm]{#1}}{\rule{0pt}{2.9cm}}}

% Captura de pantalla con su pie de figura
\newcommand{\captura}[3][width=0.92\linewidth]{%
  \begin{figure}[!htbp]\centering
  \includegraphics[#1]{evidencias/#2}
  \caption{#3}\label{fig:#2}
  \end{figure}}

\begin{document}
\sloppy

% ===========================================================================
% CARATULA
% ===========================================================================
\begin{titlepage}
\thispagestyle{empty}
\newgeometry{top=2.0cm,bottom=1.8cm,left=2.4cm,right=2.4cm}
\centering
\setlength{\parskip}{0pt}
\setstretch{1.0}

\begin{minipage}{0.48\textwidth}
\raggedright
\logo{IMAGENES/fi.png}
\end{minipage}
\begin{minipage}{0.48\textwidth}
\raggedleft
\logo{IMAGENES/unam.png}
\end{minipage}

\vspace{1.1cm}
{\Large \textbf{Universidad Nacional Autónoma de México}}\\[4pt]
{\large Facultad de Ingeniería}\\[2pt]
{\normalsize \Materia}

\vspace{1.1cm}
{\rule{\textwidth}{1.4pt}}\\[2.5pt]
{\rule{\textwidth}{0.5pt}}

\vspace{0.8cm}
{\Huge \textbf{\TituloTrabajo}}\\[12pt]
{\LARGE \SubtituloTrabajo}
\vspace{0.8cm}

{\rule{\textwidth}{0.5pt}}\\[2.5pt]
{\rule{\textwidth}{1.4pt}}

\vspace{1.3cm}

\begin{tabular}{r@{\hspace{10pt}}l}
\textbf{Semestre:} & \Semestre \\
\textbf{Profesor:} & \Profesor \\
\textbf{Grupo:} & \Grupo \\
\textbf{Fecha de entrega:} & \FechaEntrega \\
\end{tabular}

\vspace{1.1cm}

{\large \textbf{Integrantes}}\\[8pt]
{\large Basilio Illescas Andrés}\\[4pt]
{\large Cruz Macedo Samuel Santiago}\\[4pt]
{\large Fernández Hernandez Jesús Jafet}

\vfill
\restoregeometry
\end{titlepage}

\setcounter{page}{1}
{\setstretch{1.0}\setlength{\parskip}{0pt}\tableofcontents}
\newpage

% ===========================================================================
\section{Objetivo y alcance}

El programa recibe un archivo de texto con código en MiniC, un subconjunto de C, y realiza su
análisis léxico y sintáctico. El analizador léxico, generado con Flex, convierte el texto en una
secuencia de tokens con su línea y columna. El analizador sintáctico es descendente recursivo
LL(1), está escrito en C y recibe esa secuencia como único insumo para comprobar que el programa
respeta la gramática. Los errores de ambas fases se reportan con su ubicación y el análisis
continúa hasta el final del archivo. Al terminar, el programa indica si la entrada es válida.

El análisis semántico y las fases de síntesis quedan fuera del alcance, como establece la
asignación.

\noindent\textbf{Herramientas:} C (C99), Flex 2.6.4 y GNU Make.\\
\noindent\textbf{Repositorio:} \url{https://github.com/samisal123/analizador-minic}

\needspace{8\baselineskip}
\subsection{Métodos del curso aplicados}

Cada parte del proyecto sigue la técnica y la notación vistas en clase:

\begin{table}[!htbp]
\centering\small
\begin{tabularx}{\textwidth}{@{}>{\raggedright\arraybackslash}p{4.2cm}X@{}}
\toprule
\textbf{Tema} & \textbf{Dónde se aplica} \\
\midrule
2.1 a 2.3 Funciones del léxico, clases léxicas y tabla de símbolos & Tabla de tokens
(sección~\ref{sec:tokens}) y tabla de símbolos de \cmd{lexico.l} (sección~\ref{sec:lexico}) \\
2.5 y 2.6 Programación del scanner y Flex & Definiciones regulares, orden de reglas y estado
exclusivo para comentarios \\
3.1 y 3.2 Gramáticas y ambigüedad & Gramática BNF sin ambigüedad, sin recursión por la izquierda
y factorizada (sección~\ref{sec:gramatica}) \\
4.1 Conjuntos First y Follow & Sección~\ref{sec:firstfollow}, con la notación de la tarea 13 \\
4.2 Tabla de análisis y simulación con pila & Matriz $M$ con $\nexists$ en las celdas vacías y
simulaciones Pila, Cadena, Acción con \emph{pop/avanza} y \emph{fin} (secciones~\ref{sec:tabla}
y~\ref{sec:simulacion}) \\
4.3 Recuperación de errores & Modo pánico con conjunto de sincronización
(sección~\ref{sec:errores}) \\
4.4 Descendente recursivo en C & Una función por no terminal, \cmd{match()} y
\cmd{lookahead}, con los nombres \cmd{E}, \cmd{E\_Prime}, \cmd{T}, \cmd{T\_Prime} y \cmd{F}
(sección~\ref{sec:sintactico}) \\
\bottomrule
\end{tabularx}
\caption{Relación entre los temas del curso y el proyecto.}
\end{table}

% ===========================================================================
\section{Lenguaje MiniC}

\begin{itemize}[leftmargin=1.1cm,itemsep=0.25em]
\item El alfabeto lo forman las letras ASCII, los dígitos, el guion bajo y los símbolos de la
tabla~\ref{tab:tokens}. Cualquier otro carácter, incluidas las vocales acentuadas y la \cmd{ñ},
es un error léxico.
\item Se distinguen mayúsculas y minúsculas: \cmd{while} es palabra reservada y \cmd{While}, un
identificador.
\item Espacios, tabuladores y saltos de línea solo separan tokens. Hay comentarios de línea
(\cmd{//}) y de bloque (\cmd{/* */}); los de bloque pueden ocupar varias líneas y no se anidan.
\end{itemize}

\renewcommand{\arraystretch}{1.2}
\begin{table}[!htbp]
\centering\small
\begin{tabularx}{\textwidth}{@{}>{\raggedright\arraybackslash}p{3.6cm}X@{}}
\toprule
\textbf{Elemento} & \textbf{Forma} \\
\midrule
Palabras reservadas & \cmd{int}, \cmd{float}, \cmd{if}, \cmd{else}, \cmd{while} \\
Identificadores & Letra o \cmd{\_} seguida de letras, dígitos o \cmd{\_}: \cmd{x}, \cmd{\_total1} \\
Constantes & Enteras (\cmd{42}) y decimales con parte entera y fraccionaria (\cmd{3.14}) \\
Operadores & \cmd{=}; \cmd{+ - * /}; \cmd{< > <= >= == !=} \\
Delimitadores & \cmd{; , ( ) \{ \}} \\
\midrule
Declaración & \cmd{int a, b = 2;} \quad \cmd{float r = 0.5;} \\
Asignación & \cmd{id = expresión;} \\
Expresión & Operandos, paréntesis y menos unario con \cmd{+ - * /}. Precedencia: menos unario,
después \cmd{* /} y al final \cmd{+ -}; los binarios asocian por la izquierda \\
Condicional & \cmd{if (cond) \{ \dots\ \}} con \cmd{else \{ \dots\ \}} o \cmd{else if} opcional \\
Ciclo & \cmd{while (cond) \{ \dots\ \}} \\
\bottomrule
\end{tabularx}
\caption{Elementos léxicos y construcciones de MiniC.}
\label{tab:lenguaje}
\end{table}

Una condición compara dos expresiones con un operador relacional. El cuerpo de \cmd{if},
\cmd{else} y \cmd{while} siempre va entre llaves y puede estar vacío.

% ===========================================================================
\section{Tokens y expresiones regulares}
\label{sec:tokens}

Los nombres de token son los de la asignación. En la gramática cada token aparece como su
lexema, salvo \cmd{ID} (\cmd{id}), \cmd{NUMBER} (\cmd{num}) y \cmd{FIN} (\cmd{\$}).

\begin{table}[!htbp]
\centering\small
\begin{tabular}{@{}>{\ttfamily}l l @{\hspace{1.4cm}} >{\ttfamily}l l@{}}
\toprule
\normalfont\textbf{Token} & \textbf{Patrón} & \normalfont\textbf{Token} & \textbf{Patrón} \\
\midrule
INT & \cmd{int} & LESS & \cmd{<} \\
FLOAT & \cmd{float} & GREATER & \cmd{>} \\
IF & \cmd{if} & LESS\_EQUAL & \cmd{<=} \\
ELSE & \cmd{else} & GREATER\_EQUAL & \cmd{>=} \\
WHILE & \cmd{while} & EQUAL & \cmd{==} \\
ID & \cmd{\{ID\}} & NOT\_EQUAL & \cmd{!=} \\
NUMBER & \cmd{\{NUM\}} & LPAREN, RPAREN & \cmd{(} \ \cmd{)} \\
ASSIGN & \cmd{=} & LBRACE, RBRACE & \cmd{\{} \ \cmd{\}} \\
PLUS, MINUS & \cmd{+} \ \cmd{-} & SEMICOLON & \cmd{;} \\
TIMES, DIVIDE & \cmd{*} \ \cmd{/} & COMMA & \cmd{,} \\
FIN & fin de archivo & & \\
\bottomrule
\end{tabular}
\caption{Tabla de tokens.}
\label{tab:tokens}
\end{table}

\needspace{8\baselineskip}
\subsection{Definiciones regulares}

Están en la sección de definiciones de \cmd{src/lexico.l} (código~\ref{cod:lex-def}, líneas
32 a 40). \cmd{ID} y \cmd{NUM} describen los lexemas válidos. \cmd{NUM\_MAL}, \cmd{ID\_MAL} e
\cmd{ID\_UTF8} reconocen completos los lexemas mal formados (\cmd{3.14.5}, \cmd{2abc},
\cmd{año}) para reportarlos como un solo error. \cmd{UTF8} es un carácter no ASCII codificado en
UTF-8.

\begin{table}[!htbp]
\centering\small
\begin{tabular}{@{}>{\ttfamily}l l@{}}
\toprule
\normalfont\textbf{Patrón, en orden} & \textbf{Acción} \\
\midrule
{}[ \textbackslash t\textbackslash r\textbackslash n]+ \quad "//".* & Descartar \\
"/*" & Entrar al estado exclusivo \cmd{COMENTARIO} \\
"int" "float" "if" "else" "while" & Palabra reservada \\
\{ID\} \quad \{NUM\} & \cmd{ID}, \cmd{NUMBER} \\
\{NUM\_MAL\} \quad \{ID\_MAL\} \quad \{ID\_UTF8\} & Error léxico; se entrega \cmd{NUMBER} o \cmd{ID} \\
"<=" ">=" "==" "!=" & Operador relacional de dos caracteres \\
"<" ">" "=" "+" "-" "*" "/" & Operador de un carácter \\
"(" ")" "\{" "\}" ";" "," & Delimitador \\
\{UTF8\}|. & Error léxico: carácter no válido \\
\bottomrule
\end{tabular}
\caption{Reglas de la especificación Flex.}
\label{tab:reglas}
\end{table}

\needspace{8\baselineskip}
\subsection{Resolución de conflictos}

Flex elige la coincidencia más larga y, a igual longitud, la regla que aparece primero. El orden
de las reglas aprovecha ambos criterios:

\begin{itemize}[leftmargin=1.1cm,itemsep=0.2em]
\item \cmd{if} coincide con \cmd{"if"} y con \cmd{\{ID\}} con la misma longitud y gana la palabra
reservada por estar antes; en \cmd{ifx} gana \cmd{\{ID\}} por ser más larga.
\item \cmd{<=}, \cmd{//} y \cmd{/*} son más largos que \cmd{<} y \cmd{/}, así que no se parten en
dos tokens y una diagonal solo es \cmd{DIVIDE} cuando no abre un comentario.
\item En \cmd{3.14} empatan \cmd{\{NUM\}} y \cmd{\{NUM\_MAL\}} y gana \cmd{\{NUM\}}; en
\cmd{3.14.5} o \cmd{7.} solo \cmd{\{NUM\_MAL\}} cubre la cadena completa. En \cmd{2abc},
\cmd{\{ID\_MAL\}} es más larga que \cmd{\{NUM\}} y el lexema no se divide.
\end{itemize}

\needspace{8\baselineskip}
\subsection{Autómatas}

Flex traduce las expresiones a un AFD. Las figuras muestran los de identificadores y constantes,
y los de los operadores relacionales y de asignación.

\begin{figure}[!htbp]
\centering
\begin{tikzpicture}[initial text={inicio}, shorten >=1pt, node distance=2.5cm, on grid, auto,
                    >=Stealth, every state/.style={minimum size=8mm}, font=\small]
  \node[state, initial]   (q0) {$q_0$};
  \node[state, accepting] (q1) [above right=1.1cm and 3cm of q0, label=above:{\cmd{ID}}] {$q_1$};
  \node[state, accepting] (q2) [below right=1.1cm and 3cm of q0, label=below:{\cmd{NUMBER}}] {$q_2$};
  \node[state]            (q3) [right=of q2] {$q_3$};
  \node[state, accepting] (q4) [right=of q3, label=below:{\cmd{NUMBER}}] {$q_4$};
  \path[->]
    (q0) edge node {letra, \cmd{\_}} (q1)
    (q1) edge [loop right] node {letra, dígito, \cmd{\_}} ()
    (q0) edge node [swap] {dígito} (q2)
    (q2) edge [loop above] node {dígito} ()
    (q2) edge node {\cmd{.}} (q3)
    (q3) edge node {dígito} (q4)
    (q4) edge [loop above] node {dígito} ();
\end{tikzpicture}
\caption{AFD de \cmd{\{ID\}} y \cmd{\{NUM\}}. Las palabras reservadas se separan de \cmd{ID} por
la prioridad de las reglas.}
\end{figure}

\begin{figure}[!htbp]
\centering
\begin{tikzpicture}[initial text={inicio}, shorten >=1pt, node distance=2.4cm, on grid, auto,
                    >=Stealth, every state/.style={minimum size=7mm}, font=\small]
  \node[state, initial]   (q0) {$q_0$};
  \node[state, accepting] (q1) [above right=2.3cm and 2.6cm of q0, label=above:{\cmd{LESS}}] {$q_1$};
  \node[state, accepting] (q2) [right=of q1, label=above:{\cmd{LESS\_EQUAL}}] {$q_2$};
  \node[state, accepting] (q3) [above right=0.8cm and 2.6cm of q0, label=above:{\cmd{GREATER}}] {$q_3$};
  \node[state, accepting] (q4) [right=of q3, label=above:{\cmd{GREATER\_EQUAL}}] {$q_4$};
  \node[state, accepting] (q5) [below right=0.8cm and 2.6cm of q0, label=below:{\cmd{ASSIGN}}] {$q_5$};
  \node[state, accepting] (q6) [right=of q5, label=below:{\cmd{EQUAL}}] {$q_6$};
  \node[state]            (q7) [below right=2.3cm and 2.6cm of q0] {$q_7$};
  \node[state, accepting] (q8) [right=of q7, label=below:{\cmd{NOT\_EQUAL}}] {$q_8$};
  \path[->]
    (q0) edge node {\cmd{<}} (q1)  (q1) edge node {\cmd{=}} (q2)
    (q0) edge node[pos=0.6] {\cmd{>}} (q3)  (q3) edge node {\cmd{=}} (q4)
    (q0) edge node[swap,pos=0.6] {\cmd{=}} (q5)  (q5) edge node {\cmd{=}} (q6)
    (q0) edge node[swap] {\cmd{!}} (q7)  (q7) edge node {\cmd{=}} (q8);
\end{tikzpicture}
\caption{AFD de los operadores relacionales y de asignación. Un \cmd{!} sin \cmd{=} no llega a
aceptación y se reporta como carácter no válido.}
\end{figure}

% ===========================================================================
\section{Gramática}
\label{sec:gramatica}

$G = (N, T, P, S)$, con $S = \nt{programa}$, $N$ formado por los 20 no terminales que aparecen a
la izquierda de las producciones y $T$ por los terminales de la tabla~\ref{tab:tokens}. Las
producciones en BNF son:

\needspace{22\baselineskip}
\begin{lstlisting}
<programa>       ::= <lista_sent>
<lista_sent>     ::= <sentencia> <lista_sent> | ε
<sentencia>      ::= <declaracion> | <asignacion> | <sent_if> | <sent_while>
<declaracion>    ::= <tipo> id <inicializacion> <lista_id> ;
<tipo>           ::= int | float
<lista_id>       ::= , id <inicializacion> <lista_id> | ε
<inicializacion> ::= = <E> | ε
<asignacion>     ::= id = <E> ;
<sent_if>        ::= if ( <condicion> ) <bloque> <parte_else>
<parte_else>     ::= else <resto_else> | ε
<resto_else>     ::= <bloque> | <sent_if>
<sent_while>     ::= while ( <condicion> ) <bloque>
<bloque>         ::= { <lista_sent> }
<condicion>      ::= <E> <op_rel> <E>
<op_rel>         ::= < | > | <= | >= | == | !=
<E>              ::= <T> <E'>
<E'>             ::= + <T> <E'> | - <T> <E'> | ε
<T>              ::= <F> <T'>
<T'>             ::= * <F> <T'> | / <F> <T'> | ε
<F>              ::= ( <E> ) | id | num | - <F>
\end{lstlisting}

Para la tabla LL(1) cada alternativa se numera por separado, como en clase:

\begin{table}[!htbp]
\centering\footnotesize
\renewcommand{\arraystretch}{1.1}
\begin{tabular}{@{}r l @{\hspace{0.8cm}} r l@{}}
\toprule
1 & \nt{programa} $\rightarrow$ \nt{lista\_sent} & 22 & \nt{bloque} $\rightarrow$ \texttt{\{} \nt{lista\_sent} \texttt{\}} \\
2 & \nt{lista\_sent} $\rightarrow$ \nt{sentencia} \nt{lista\_sent} & 23 & \nt{condicion} $\rightarrow$ \nt{E} \nt{op\_rel} \nt{E} \\
3 & \nt{lista\_sent} $\rightarrow$ $\varepsilon$ & 24 & \nt{op\_rel} $\rightarrow$ \texttt{<} \\
4 & \nt{sentencia} $\rightarrow$ \nt{declaracion} & 25 & \nt{op\_rel} $\rightarrow$ \texttt{>} \\
5 & \nt{sentencia} $\rightarrow$ \nt{asignacion} & 26 & \nt{op\_rel} $\rightarrow$ \texttt{<=} \\
6 & \nt{sentencia} $\rightarrow$ \nt{sent\_if} & 27 & \nt{op\_rel} $\rightarrow$ \texttt{>=} \\
7 & \nt{sentencia} $\rightarrow$ \nt{sent\_while} & 28 & \nt{op\_rel} $\rightarrow$ \texttt{==} \\
8 & \nt{declaracion} $\rightarrow$ \nt{tipo} \texttt{id} \nt{inicializacion} \nt{lista\_id} \texttt{;} & 29 & \nt{op\_rel} $\rightarrow$ \texttt{!=} \\
9 & \nt{tipo} $\rightarrow$ \texttt{int} & 30 & \nt{E} $\rightarrow$ \nt{T} \nt{E'} \\
10 & \nt{tipo} $\rightarrow$ \texttt{float} & 31 & \nt{E'} $\rightarrow$ \texttt{+} \nt{T} \nt{E'} \\
11 & \nt{lista\_id} $\rightarrow$ \texttt{,} \texttt{id} \nt{inicializacion} \nt{lista\_id} & 32 & \nt{E'} $\rightarrow$ \texttt{-} \nt{T} \nt{E'} \\
12 & \nt{lista\_id} $\rightarrow$ $\varepsilon$ & 33 & \nt{E'} $\rightarrow$ $\varepsilon$ \\
13 & \nt{inicializacion} $\rightarrow$ \texttt{=} \nt{E} & 34 & \nt{T} $\rightarrow$ \nt{F} \nt{T'} \\
14 & \nt{inicializacion} $\rightarrow$ $\varepsilon$ & 35 & \nt{T'} $\rightarrow$ \texttt{*} \nt{F} \nt{T'} \\
15 & \nt{asignacion} $\rightarrow$ \texttt{id} \texttt{=} \nt{E} \texttt{;} & 36 & \nt{T'} $\rightarrow$ \texttt{/} \nt{F} \nt{T'} \\
16 & \nt{sent\_if} $\rightarrow$ \texttt{if} \texttt{(} \nt{condicion} \texttt{)} \nt{bloque} \nt{parte\_else} & 37 & \nt{T'} $\rightarrow$ $\varepsilon$ \\
17 & \nt{parte\_else} $\rightarrow$ \texttt{else} \nt{resto\_else} & 38 & \nt{F} $\rightarrow$ \texttt{(} \nt{E} \texttt{)} \\
18 & \nt{parte\_else} $\rightarrow$ $\varepsilon$ & 39 & \nt{F} $\rightarrow$ \texttt{id} \\
19 & \nt{resto\_else} $\rightarrow$ \nt{bloque} & 40 & \nt{F} $\rightarrow$ \texttt{num} \\
20 & \nt{resto\_else} $\rightarrow$ \nt{sent\_if} & 41 & \nt{F} $\rightarrow$ \texttt{-} \nt{F} \\
21 & \nt{sent\_while} $\rightarrow$ \texttt{while} \texttt{(} \nt{condicion} \texttt{)} \nt{bloque} & & \\
\bottomrule
\end{tabular}
\caption{Producciones numeradas.}
\label{tab:producciones}
\end{table}

\needspace{8\baselineskip}
\subsection{Tratamiento de la ambigüedad}

\begin{itemize}[leftmargin=1.1cm,itemsep=0.3em]
\item \textbf{Precedencia.} $E \rightarrow E + E \mid E * E \mid \mathrm{id}$ da dos árboles para
\cmd{id + id * id}. Aquí cada nivel solo contiene al siguiente: \nt{E} suma términos, \nt{T}
multiplica factores y \nt{F} tiene los operandos, los paréntesis y el menos unario. Así \cmd{*}
y \cmd{/} se agrupan antes que \cmd{+} y \cmd{-}, y cada cadena tiene un solo árbol.
\item \textbf{Recursión por la izquierda.} $E \rightarrow E + T \mid T$ haría que \cmd{E()} se
llamara a sí misma sin consumir tokens. Con $A \rightarrow A\alpha \mid \beta \Rightarrow A
\rightarrow \beta A',\ A' \rightarrow \alpha A' \mid \varepsilon$ quedan \nt{E} y \nt{E'}; lo
mismo para \nt{T}. Es la misma gramática de expresiones de la tarea 13.
\item \textbf{Factorización.} El \cmd{if} con y sin \cmd{else} comparte el prefijo
\cmd{if ( condicion ) bloque}, que queda una sola vez en \nt{sent\_if}; la decisión pasa a
\nt{parte\_else}. \cmd{else \{\}} y \cmd{else if} comparten \cmd{else} y se separan en
\nt{resto\_else}.
\item \textbf{\cmd{else} colgante.} Como el cuerpo de un \cmd{if} siempre es un bloque entre
llaves, cada \cmd{else} pertenece a un solo \cmd{if}. En la gramática, \cmd{else} no está en
Follow(\nt{parte\_else}) y la celda $M[\nt{parte\_else}, \texttt{else}]$ tiene una sola
producción.
\end{itemize}

% ===========================================================================
\section{Conjuntos First y Follow}
\label{sec:firstfollow}

Se calcularon con las reglas del tema 4.1, iterando hasta que ningún conjunto cambió. El símbolo
inicial lleva \cmd{\$} en su Follow.

\begin{table}[!htbp]
\centering\footnotesize
\renewcommand{\arraystretch}{1.25}
\begin{tabular}{@{}l >{\raggedright\arraybackslash}p{5.5cm} >{\raggedright\arraybackslash}p{6.3cm}@{}}
\toprule
\textbf{No terminal} & \textbf{First} & \textbf{Follow} \\
\midrule
\nt{programa} & \{ \texttt{int}, \texttt{float}, \texttt{if}, \texttt{while}, \texttt{id}, $\varepsilon$ \} & \{ \texttt{\$} \} \\
\nt{lista\_sent} & \{ \texttt{int}, \texttt{float}, \texttt{if}, \texttt{while}, \texttt{id}, $\varepsilon$ \} & \{ \texttt{\}}, \texttt{\$} \} \\
\nt{sentencia} & \{ \texttt{int}, \texttt{float}, \texttt{if}, \texttt{while}, \texttt{id} \} & \{ \texttt{int}, \texttt{float}, \texttt{if}, \texttt{while}, \texttt{id}, \texttt{\}}, \texttt{\$} \} \\
\nt{declaracion} & \{ \texttt{int}, \texttt{float} \} & \{ \texttt{int}, \texttt{float}, \texttt{if}, \texttt{while}, \texttt{id}, \texttt{\}}, \texttt{\$} \} \\
\nt{tipo} & \{ \texttt{int}, \texttt{float} \} & \{ \texttt{id} \} \\
\nt{lista\_id} & \{ \texttt{,}, $\varepsilon$ \} & \{ \texttt{;} \} \\
\nt{inicializacion} & \{ \texttt{=}, $\varepsilon$ \} & \{ \texttt{;}, \texttt{,} \} \\
\nt{asignacion} & \{ \texttt{id} \} & \{ \texttt{int}, \texttt{float}, \texttt{if}, \texttt{while}, \texttt{id}, \texttt{\}}, \texttt{\$} \} \\
\nt{sent\_if} & \{ \texttt{if} \} & \{ \texttt{int}, \texttt{float}, \texttt{if}, \texttt{while}, \texttt{id}, \texttt{\}}, \texttt{\$} \} \\
\nt{parte\_else} & \{ \texttt{else}, $\varepsilon$ \} & \{ \texttt{int}, \texttt{float}, \texttt{if}, \texttt{while}, \texttt{id}, \texttt{\}}, \texttt{\$} \} \\
\nt{resto\_else} & \{ \texttt{if}, \texttt{\{} \} & \{ \texttt{int}, \texttt{float}, \texttt{if}, \texttt{while}, \texttt{id}, \texttt{\}}, \texttt{\$} \} \\
\nt{sent\_while} & \{ \texttt{while} \} & \{ \texttt{int}, \texttt{float}, \texttt{if}, \texttt{while}, \texttt{id}, \texttt{\}}, \texttt{\$} \} \\
\nt{bloque} & \{ \texttt{\{} \} & \{ \texttt{int}, \texttt{float}, \texttt{if}, \texttt{else}, \texttt{while}, \texttt{id}, \texttt{\}}, \texttt{\$} \} \\
\nt{condicion} & \{ \texttt{id}, \texttt{num}, \texttt{-}, \texttt{(} \} & \{ \texttt{)} \} \\
\nt{op\_rel} & \{ \texttt{<}, \texttt{>}, \texttt{<=}, \texttt{>=}, \texttt{==}, \texttt{!=} \} & \{ \texttt{id}, \texttt{num}, \texttt{-}, \texttt{(} \} \\
\nt{E} & \{ \texttt{id}, \texttt{num}, \texttt{-}, \texttt{(} \} & \{ \texttt{<}, \texttt{>}, \texttt{<=}, \texttt{>=}, \texttt{==}, \texttt{!=}, \texttt{)}, \texttt{;}, \texttt{,} \} \\
\nt{E'} & \{ \texttt{+}, \texttt{-}, $\varepsilon$ \} & \{ \texttt{<}, \texttt{>}, \texttt{<=}, \texttt{>=}, \texttt{==}, \texttt{!=}, \texttt{)}, \texttt{;}, \texttt{,} \} \\
\nt{T} & \{ \texttt{id}, \texttt{num}, \texttt{-}, \texttt{(} \} & \{ \texttt{+}, \texttt{-}, \texttt{<}, \texttt{>}, \texttt{<=}, \texttt{>=}, \texttt{==}, \texttt{!=}, \texttt{)}, \texttt{;}, \texttt{,} \} \\
\nt{T'} & \{ \texttt{*}, \texttt{/}, $\varepsilon$ \} & \{ \texttt{+}, \texttt{-}, \texttt{<}, \texttt{>}, \texttt{<=}, \texttt{>=}, \texttt{==}, \texttt{!=}, \texttt{)}, \texttt{;}, \texttt{,} \} \\
\nt{F} & \{ \texttt{id}, \texttt{num}, \texttt{-}, \texttt{(} \} & \{ \texttt{+}, \texttt{-}, \texttt{*}, \texttt{/}, \texttt{<}, \texttt{>}, \texttt{<=}, \texttt{>=}, \texttt{==}, \texttt{!=}, \texttt{)}, \texttt{;}, \texttt{,} \} \\
\bottomrule
\end{tabular}
\caption{Conjuntos First y Follow.}
\label{tab:firstfollow}
\end{table}

% ===========================================================================
\section{Tabla de análisis LL(1)}
\label{sec:tabla}

Para cada producción $A \rightarrow \alpha$ se escribe su número en $M[A, a]$ para todo
$a \in \text{First}(\alpha)$ y, si $\varepsilon \in \text{First}(\alpha)$, también en
$M[A, b]$ para todo $b \in \text{Follow}(A)$. Las celdas sin producción llevan $\nexists$ (no
hay correspondencia), igual que en la tarea 13. Por su ancho, la matriz se divide en dos partes (tablas~\ref{tab:matriz} y~\ref{tab:matriz2}) y usa la numeración de la tabla~\ref{tab:producciones}. Ninguna celda tiene más de una producción, así que la gramática
es LL(1).

\begin{table}[p]
\centering\footnotesize
\setlength{\tabcolsep}{5pt}
\renewcommand{\arraystretch}{1.08}
\begin{tabular}{@{}l*{13}{c}@{}}
\toprule
\textbf{No terminal} & \textbf{\texttt{int}} & \textbf{\texttt{float}} & \textbf{\texttt{if}} & \textbf{\texttt{else}} & \textbf{\texttt{while}} & \textbf{\texttt{id}} & \textbf{\texttt{num}} & \textbf{\texttt{=}} & \textbf{\texttt{+}} & \textbf{\texttt{-}} & \textbf{\texttt{*}} & \textbf{\texttt{/}} & \textbf{\texttt{<}} \\
\midrule
\nt{programa} & 1 & 1 & 1 & $\nexists$ & 1 & 1 & $\nexists$ & $\nexists$ & $\nexists$ & $\nexists$ & $\nexists$ & $\nexists$ & $\nexists$ \\
\nt{lista\_sent} & 2 & 2 & 2 & $\nexists$ & 2 & 2 & $\nexists$ & $\nexists$ & $\nexists$ & $\nexists$ & $\nexists$ & $\nexists$ & $\nexists$ \\
\nt{sentencia} & 4 & 4 & 6 & $\nexists$ & 7 & 5 & $\nexists$ & $\nexists$ & $\nexists$ & $\nexists$ & $\nexists$ & $\nexists$ & $\nexists$ \\
\nt{declaracion} & 8 & 8 & $\nexists$ & $\nexists$ & $\nexists$ & $\nexists$ & $\nexists$ & $\nexists$ & $\nexists$ & $\nexists$ & $\nexists$ & $\nexists$ & $\nexists$ \\
\nt{tipo} & 9 & 10 & $\nexists$ & $\nexists$ & $\nexists$ & $\nexists$ & $\nexists$ & $\nexists$ & $\nexists$ & $\nexists$ & $\nexists$ & $\nexists$ & $\nexists$ \\
\nt{lista\_id} & $\nexists$ & $\nexists$ & $\nexists$ & $\nexists$ & $\nexists$ & $\nexists$ & $\nexists$ & $\nexists$ & $\nexists$ & $\nexists$ & $\nexists$ & $\nexists$ & $\nexists$ \\
\nt{inicializacion} & $\nexists$ & $\nexists$ & $\nexists$ & $\nexists$ & $\nexists$ & $\nexists$ & $\nexists$ & 13 & $\nexists$ & $\nexists$ & $\nexists$ & $\nexists$ & $\nexists$ \\
\nt{asignacion} & $\nexists$ & $\nexists$ & $\nexists$ & $\nexists$ & $\nexists$ & 15 & $\nexists$ & $\nexists$ & $\nexists$ & $\nexists$ & $\nexists$ & $\nexists$ & $\nexists$ \\
\nt{sent\_if} & $\nexists$ & $\nexists$ & 16 & $\nexists$ & $\nexists$ & $\nexists$ & $\nexists$ & $\nexists$ & $\nexists$ & $\nexists$ & $\nexists$ & $\nexists$ & $\nexists$ \\
\nt{parte\_else} & 18 & 18 & 18 & 17 & 18 & 18 & $\nexists$ & $\nexists$ & $\nexists$ & $\nexists$ & $\nexists$ & $\nexists$ & $\nexists$ \\
\nt{resto\_else} & $\nexists$ & $\nexists$ & 20 & $\nexists$ & $\nexists$ & $\nexists$ & $\nexists$ & $\nexists$ & $\nexists$ & $\nexists$ & $\nexists$ & $\nexists$ & $\nexists$ \\
\nt{sent\_while} & $\nexists$ & $\nexists$ & $\nexists$ & $\nexists$ & 21 & $\nexists$ & $\nexists$ & $\nexists$ & $\nexists$ & $\nexists$ & $\nexists$ & $\nexists$ & $\nexists$ \\
\nt{bloque} & $\nexists$ & $\nexists$ & $\nexists$ & $\nexists$ & $\nexists$ & $\nexists$ & $\nexists$ & $\nexists$ & $\nexists$ & $\nexists$ & $\nexists$ & $\nexists$ & $\nexists$ \\
\nt{condicion} & $\nexists$ & $\nexists$ & $\nexists$ & $\nexists$ & $\nexists$ & 23 & 23 & $\nexists$ & $\nexists$ & 23 & $\nexists$ & $\nexists$ & $\nexists$ \\
\nt{op\_rel} & $\nexists$ & $\nexists$ & $\nexists$ & $\nexists$ & $\nexists$ & $\nexists$ & $\nexists$ & $\nexists$ & $\nexists$ & $\nexists$ & $\nexists$ & $\nexists$ & 24 \\
\nt{E} & $\nexists$ & $\nexists$ & $\nexists$ & $\nexists$ & $\nexists$ & 30 & 30 & $\nexists$ & $\nexists$ & 30 & $\nexists$ & $\nexists$ & $\nexists$ \\
\nt{E'} & $\nexists$ & $\nexists$ & $\nexists$ & $\nexists$ & $\nexists$ & $\nexists$ & $\nexists$ & $\nexists$ & 31 & 32 & $\nexists$ & $\nexists$ & 33 \\
\nt{T} & $\nexists$ & $\nexists$ & $\nexists$ & $\nexists$ & $\nexists$ & 34 & 34 & $\nexists$ & $\nexists$ & 34 & $\nexists$ & $\nexists$ & $\nexists$ \\
\nt{T'} & $\nexists$ & $\nexists$ & $\nexists$ & $\nexists$ & $\nexists$ & $\nexists$ & $\nexists$ & $\nexists$ & 37 & 37 & 35 & 36 & 37 \\
\nt{F} & $\nexists$ & $\nexists$ & $\nexists$ & $\nexists$ & $\nexists$ & 39 & 40 & $\nexists$ & $\nexists$ & 41 & $\nexists$ & $\nexists$ & $\nexists$ \\
\bottomrule
\end{tabular}
\caption{Tabla de análisis LL(1), columnas \texttt{int} a \texttt{<}. Cada número es una producción de la tabla~\ref{tab:producciones}; $\nexists$ indica error.}
\label{tab:matriz}

\vspace{1.2em}
\footnotesize
\setlength{\tabcolsep}{5pt}
\renewcommand{\arraystretch}{1.08}
\begin{tabular}{@{}l*{12}{c}@{}}
\toprule
\textbf{No terminal} & \textbf{\texttt{>}} & \textbf{\texttt{<=}} & \textbf{\texttt{>=}} & \textbf{\texttt{==}} & \textbf{\texttt{!=}} & \textbf{\texttt{(}} & \textbf{\texttt{)}} & \textbf{\texttt{\{}} & \textbf{\texttt{\}}} & \textbf{\texttt{;}} & \textbf{\texttt{,}} & \textbf{\texttt{\$}} \\
\midrule
\nt{programa} & $\nexists$ & $\nexists$ & $\nexists$ & $\nexists$ & $\nexists$ & $\nexists$ & $\nexists$ & $\nexists$ & $\nexists$ & $\nexists$ & $\nexists$ & 1 \\
\nt{lista\_sent} & $\nexists$ & $\nexists$ & $\nexists$ & $\nexists$ & $\nexists$ & $\nexists$ & $\nexists$ & $\nexists$ & 3 & $\nexists$ & $\nexists$ & 3 \\
\nt{sentencia} & $\nexists$ & $\nexists$ & $\nexists$ & $\nexists$ & $\nexists$ & $\nexists$ & $\nexists$ & $\nexists$ & $\nexists$ & $\nexists$ & $\nexists$ & $\nexists$ \\
\nt{declaracion} & $\nexists$ & $\nexists$ & $\nexists$ & $\nexists$ & $\nexists$ & $\nexists$ & $\nexists$ & $\nexists$ & $\nexists$ & $\nexists$ & $\nexists$ & $\nexists$ \\
\nt{tipo} & $\nexists$ & $\nexists$ & $\nexists$ & $\nexists$ & $\nexists$ & $\nexists$ & $\nexists$ & $\nexists$ & $\nexists$ & $\nexists$ & $\nexists$ & $\nexists$ \\
\nt{lista\_id} & $\nexists$ & $\nexists$ & $\nexists$ & $\nexists$ & $\nexists$ & $\nexists$ & $\nexists$ & $\nexists$ & $\nexists$ & 12 & 11 & $\nexists$ \\
\nt{inicializacion} & $\nexists$ & $\nexists$ & $\nexists$ & $\nexists$ & $\nexists$ & $\nexists$ & $\nexists$ & $\nexists$ & $\nexists$ & 14 & 14 & $\nexists$ \\
\nt{asignacion} & $\nexists$ & $\nexists$ & $\nexists$ & $\nexists$ & $\nexists$ & $\nexists$ & $\nexists$ & $\nexists$ & $\nexists$ & $\nexists$ & $\nexists$ & $\nexists$ \\
\nt{sent\_if} & $\nexists$ & $\nexists$ & $\nexists$ & $\nexists$ & $\nexists$ & $\nexists$ & $\nexists$ & $\nexists$ & $\nexists$ & $\nexists$ & $\nexists$ & $\nexists$ \\
\nt{parte\_else} & $\nexists$ & $\nexists$ & $\nexists$ & $\nexists$ & $\nexists$ & $\nexists$ & $\nexists$ & $\nexists$ & 18 & $\nexists$ & $\nexists$ & 18 \\
\nt{resto\_else} & $\nexists$ & $\nexists$ & $\nexists$ & $\nexists$ & $\nexists$ & $\nexists$ & $\nexists$ & 19 & $\nexists$ & $\nexists$ & $\nexists$ & $\nexists$ \\
\nt{sent\_while} & $\nexists$ & $\nexists$ & $\nexists$ & $\nexists$ & $\nexists$ & $\nexists$ & $\nexists$ & $\nexists$ & $\nexists$ & $\nexists$ & $\nexists$ & $\nexists$ \\
\nt{bloque} & $\nexists$ & $\nexists$ & $\nexists$ & $\nexists$ & $\nexists$ & $\nexists$ & $\nexists$ & 22 & $\nexists$ & $\nexists$ & $\nexists$ & $\nexists$ \\
\nt{condicion} & $\nexists$ & $\nexists$ & $\nexists$ & $\nexists$ & $\nexists$ & 23 & $\nexists$ & $\nexists$ & $\nexists$ & $\nexists$ & $\nexists$ & $\nexists$ \\
\nt{op\_rel} & 25 & 26 & 27 & 28 & 29 & $\nexists$ & $\nexists$ & $\nexists$ & $\nexists$ & $\nexists$ & $\nexists$ & $\nexists$ \\
\nt{E} & $\nexists$ & $\nexists$ & $\nexists$ & $\nexists$ & $\nexists$ & 30 & $\nexists$ & $\nexists$ & $\nexists$ & $\nexists$ & $\nexists$ & $\nexists$ \\
\nt{E'} & 33 & 33 & 33 & 33 & 33 & $\nexists$ & 33 & $\nexists$ & $\nexists$ & 33 & 33 & $\nexists$ \\
\nt{T} & $\nexists$ & $\nexists$ & $\nexists$ & $\nexists$ & $\nexists$ & 34 & $\nexists$ & $\nexists$ & $\nexists$ & $\nexists$ & $\nexists$ & $\nexists$ \\
\nt{T'} & 37 & 37 & 37 & 37 & 37 & $\nexists$ & 37 & $\nexists$ & $\nexists$ & 37 & 37 & $\nexists$ \\
\nt{F} & $\nexists$ & $\nexists$ & $\nexists$ & $\nexists$ & $\nexists$ & 38 & $\nexists$ & $\nexists$ & $\nexists$ & $\nexists$ & $\nexists$ & $\nexists$ \\
\bottomrule
\end{tabular}
\caption{Tabla de análisis LL(1), columnas \texttt{>} a \texttt{\$}.}
\label{tab:matriz2}
\end{table}


% ===========================================================================
\section{Simulación del análisis}
\label{sec:simulacion}

Las simulaciones siguen el formato de clase: la pila empieza con el símbolo inicial y
\cmd{\$}, con el tope a la izquierda. Si el tope es no terminal se aplica la producción de
$M[\text{tope}, \text{primer símbolo de la cadena}]$; si es terminal y coincide con la cadena, se
hace \emph{pop/avanza}; con \cmd{\$} en ambas columnas se termina con \emph{fin}.

\needspace{8\baselineskip}
\subsection{Simulación 1: \texttt{y = x + 5;} (válida)}

El analizador léxico entrega \cmd{id = id + num ;}.

\begin{table}[!htbp]
\centering\footnotesize
\renewcommand{\arraystretch}{1.08}
\begin{tabular}{@{}>{\raggedright\arraybackslash}p{5.0cm} >{\raggedright\arraybackslash}p{3.7cm} >{\raggedright\arraybackslash}p{5.6cm}@{}}
\toprule
\textbf{Pila} & \textbf{Cadena} & \textbf{Acción} \\
\midrule
\nt{programa} \texttt{\$} & \texttt{id} \texttt{=} \texttt{id} \texttt{+} \texttt{num} \texttt{;} \texttt{\$} & \nt{programa} $\rightarrow$ \nt{lista\_sent} \\
\nt{lista\_sent} \texttt{\$} & \texttt{id} \texttt{=} \texttt{id} \texttt{+} \texttt{num} \texttt{;} \texttt{\$} & \nt{lista\_sent} $\rightarrow$ \nt{sentencia} \nt{lista\_sent} \\
\nt{sentencia} \nt{lista\_sent} \texttt{\$} & \texttt{id} \texttt{=} \texttt{id} \texttt{+} \texttt{num} \texttt{;} \texttt{\$} & \nt{sentencia} $\rightarrow$ \nt{asignacion} \\
\nt{asignacion} \nt{lista\_sent} \texttt{\$} & \texttt{id} \texttt{=} \texttt{id} \texttt{+} \texttt{num} \texttt{;} \texttt{\$} & \nt{asignacion} $\rightarrow$ \texttt{id} \texttt{=} \nt{E} \texttt{;} \\
\texttt{id} \texttt{=} \nt{E} \texttt{;} \nt{lista\_sent} \texttt{\$} & \texttt{id} \texttt{=} \texttt{id} \texttt{+} \texttt{num} \texttt{;} \texttt{\$} & pop/avanza \\
\texttt{=} \nt{E} \texttt{;} \nt{lista\_sent} \texttt{\$} & \texttt{=} \texttt{id} \texttt{+} \texttt{num} \texttt{;} \texttt{\$} & pop/avanza \\
\nt{E} \texttt{;} \nt{lista\_sent} \texttt{\$} & \texttt{id} \texttt{+} \texttt{num} \texttt{;} \texttt{\$} & \nt{E} $\rightarrow$ \nt{T} \nt{E'} \\
\nt{T} \nt{E'} \texttt{;} \nt{lista\_sent} \texttt{\$} & \texttt{id} \texttt{+} \texttt{num} \texttt{;} \texttt{\$} & \nt{T} $\rightarrow$ \nt{F} \nt{T'} \\
\nt{F} \nt{T'} \nt{E'} \texttt{;} \nt{lista\_sent} \texttt{\$} & \texttt{id} \texttt{+} \texttt{num} \texttt{;} \texttt{\$} & \nt{F} $\rightarrow$ \texttt{id} \\
\texttt{id} \nt{T'} \nt{E'} \texttt{;} \nt{lista\_sent} \texttt{\$} & \texttt{id} \texttt{+} \texttt{num} \texttt{;} \texttt{\$} & pop/avanza \\
\nt{T'} \nt{E'} \texttt{;} \nt{lista\_sent} \texttt{\$} & \texttt{+} \texttt{num} \texttt{;} \texttt{\$} & \nt{T'} $\rightarrow$ $\varepsilon$ \\
\nt{E'} \texttt{;} \nt{lista\_sent} \texttt{\$} & \texttt{+} \texttt{num} \texttt{;} \texttt{\$} & \nt{E'} $\rightarrow$ \texttt{+} \nt{T} \nt{E'} \\
\texttt{+} \nt{T} \nt{E'} \texttt{;} \nt{lista\_sent} \texttt{\$} & \texttt{+} \texttt{num} \texttt{;} \texttt{\$} & pop/avanza \\
\nt{T} \nt{E'} \texttt{;} \nt{lista\_sent} \texttt{\$} & \texttt{num} \texttt{;} \texttt{\$} & \nt{T} $\rightarrow$ \nt{F} \nt{T'} \\
\nt{F} \nt{T'} \nt{E'} \texttt{;} \nt{lista\_sent} \texttt{\$} & \texttt{num} \texttt{;} \texttt{\$} & \nt{F} $\rightarrow$ \texttt{num} \\
\texttt{num} \nt{T'} \nt{E'} \texttt{;} \nt{lista\_sent} \texttt{\$} & \texttt{num} \texttt{;} \texttt{\$} & pop/avanza \\
\nt{T'} \nt{E'} \texttt{;} \nt{lista\_sent} \texttt{\$} & \texttt{;} \texttt{\$} & \nt{T'} $\rightarrow$ $\varepsilon$ \\
\nt{E'} \texttt{;} \nt{lista\_sent} \texttt{\$} & \texttt{;} \texttt{\$} & \nt{E'} $\rightarrow$ $\varepsilon$ \\
\texttt{;} \nt{lista\_sent} \texttt{\$} & \texttt{;} \texttt{\$} & pop/avanza \\
\nt{lista\_sent} \texttt{\$} & \texttt{\$} & \nt{lista\_sent} $\rightarrow$ $\varepsilon$ \\
\texttt{\$} & \texttt{\$} & fin \\
\bottomrule
\end{tabular}
\caption{Simulación de \cmd{y = x + 5;}.}
\label{tab:sim1}
\end{table}

\begin{figure}[!htbp]
\centering
\begin{tikzpicture}[level distance=8.5mm, font=\small,
  edge from parent/.style={draw, line width=0.4pt},
  level 1/.style={sibling distance=38mm}, level 2/.style={sibling distance=38mm},
  level 3/.style={sibling distance=19mm}, level 4/.style={sibling distance=19mm},
  level 5/.style={sibling distance=34mm}, level 6/.style={sibling distance=13mm},
  level 7/.style={sibling distance=11mm}]
\node {\nt{programa}}
  child {node {\nt{lista\_sent}}
    child {node {\nt{sentencia}}
      child {node {\nt{asignacion}}
        child {node {\cmd{id} (y)}}
        child {node {\cmd{=}}}
        child {node {\nt{E}}
          child {node {\nt{T}}
            child {node {\nt{F}} child {node {\cmd{id} (x)}}}
            child {node {\nt{T'}} child {node {\eps}}}}
          child {node {\nt{E'}}
            child {node {\cmd{+}}}
            child {node {\nt{T}}
              child {node {\nt{F}} child {node {\cmd{num} (5)}}}
              child {node {\nt{T'}} child {node {\eps}}}}
            child {node {\nt{E'}} child {node {\eps}}}}}
        child {node {\cmd{;}}}}}
    child {node {\nt{lista\_sent}} child {node {\eps}}}};
\end{tikzpicture}
\caption{Árbol sintáctico de \cmd{y = x + 5;}. Cada nodo interno es una llamada a función del
analizador y cada hoja, un \cmd{match()}.}
\label{fig:arbol}
\end{figure}

\needspace{8\baselineskip}
\subsection{Simulación 2: \texttt{x = (10 + 5;} (inválida)}

Es el caso de error de la asignación. La cadena de tokens es \cmd{id = ( num + num ;}.

\begin{table}[!htbp]
\centering\footnotesize
\renewcommand{\arraystretch}{1.08}
\begin{tabular}{@{}>{\raggedright\arraybackslash}p{5.0cm} >{\raggedright\arraybackslash}p{3.7cm} >{\raggedright\arraybackslash}p{5.6cm}@{}}
\toprule
\textbf{Pila} & \textbf{Cadena} & \textbf{Acción} \\
\midrule
\nt{programa} \texttt{\$} & \texttt{id} \texttt{=} \texttt{(} \texttt{num} \texttt{+} \texttt{num} \texttt{;} \texttt{\$} & \nt{programa} $\rightarrow$ \nt{lista\_sent} \\
\nt{lista\_sent} \texttt{\$} & \texttt{id} \texttt{=} \texttt{(} \texttt{num} \texttt{+} \texttt{num} \texttt{;} \texttt{\$} & \nt{lista\_sent} $\rightarrow$ \nt{sentencia} \nt{lista\_sent} \\
\nt{sentencia} \nt{lista\_sent} \texttt{\$} & \texttt{id} \texttt{=} \texttt{(} \texttt{num} \texttt{+} \texttt{num} \texttt{;} \texttt{\$} & \nt{sentencia} $\rightarrow$ \nt{asignacion} \\
\nt{asignacion} \nt{lista\_sent} \texttt{\$} & \texttt{id} \texttt{=} \texttt{(} \texttt{num} \texttt{+} \texttt{num} \texttt{;} \texttt{\$} & \nt{asignacion} $\rightarrow$ \texttt{id} \texttt{=} \nt{E} \texttt{;} \\
\texttt{id} \texttt{=} \nt{E} \texttt{;} \nt{lista\_sent} \texttt{\$} & \texttt{id} \texttt{=} \texttt{(} \texttt{num} \texttt{+} \texttt{num} \texttt{;} \texttt{\$} & pop/avanza \\
\texttt{=} \nt{E} \texttt{;} \nt{lista\_sent} \texttt{\$} & \texttt{=} \texttt{(} \texttt{num} \texttt{+} \texttt{num} \texttt{;} \texttt{\$} & pop/avanza \\
\nt{E} \texttt{;} \nt{lista\_sent} \texttt{\$} & \texttt{(} \texttt{num} \texttt{+} \texttt{num} \texttt{;} \texttt{\$} & \nt{E} $\rightarrow$ \nt{T} \nt{E'} \\
\nt{T} \nt{E'} \texttt{;} \nt{lista\_sent} \texttt{\$} & \texttt{(} \texttt{num} \texttt{+} \texttt{num} \texttt{;} \texttt{\$} & \nt{T} $\rightarrow$ \nt{F} \nt{T'} \\
\nt{F} \nt{T'} \nt{E'} \texttt{;} \nt{lista\_sent} \texttt{\$} & \texttt{(} \texttt{num} \texttt{+} \texttt{num} \texttt{;} \texttt{\$} & \nt{F} $\rightarrow$ \texttt{(} \nt{E} \texttt{)} \\
\texttt{(} \nt{E} \texttt{)} \nt{T'} \nt{E'} \texttt{;} \nt{lista\_sent} \texttt{\$} & \texttt{(} \texttt{num} \texttt{+} \texttt{num} \texttt{;} \texttt{\$} & pop/avanza \\
\nt{E} \texttt{)} \nt{T'} \nt{E'} \texttt{;} \nt{lista\_sent} \texttt{\$} & \texttt{num} \texttt{+} \texttt{num} \texttt{;} \texttt{\$} & \nt{E} $\rightarrow$ \nt{T} \nt{E'} \\
\nt{T} \nt{E'} \texttt{)} \nt{T'} \nt{E'} \texttt{;} \nt{lista\_sent} \texttt{\$} & \texttt{num} \texttt{+} \texttt{num} \texttt{;} \texttt{\$} & \nt{T} $\rightarrow$ \nt{F} \nt{T'} \\
\nt{F} \nt{T'} \nt{E'} \texttt{)} \nt{T'} \nt{E'} \texttt{;} \nt{lista\_sent} \texttt{\$} & \texttt{num} \texttt{+} \texttt{num} \texttt{;} \texttt{\$} & \nt{F} $\rightarrow$ \texttt{num} \\
\texttt{num} \nt{T'} \nt{E'} \texttt{)} \nt{T'} \nt{E'} \texttt{;} \nt{lista\_sent} \texttt{\$} & \texttt{num} \texttt{+} \texttt{num} \texttt{;} \texttt{\$} & pop/avanza \\
\nt{T'} \nt{E'} \texttt{)} \nt{T'} \nt{E'} \texttt{;} \nt{lista\_sent} \texttt{\$} & \texttt{+} \texttt{num} \texttt{;} \texttt{\$} & \nt{T'} $\rightarrow$ $\varepsilon$ \\
\nt{E'} \texttt{)} \nt{T'} \nt{E'} \texttt{;} \nt{lista\_sent} \texttt{\$} & \texttt{+} \texttt{num} \texttt{;} \texttt{\$} & \nt{E'} $\rightarrow$ \texttt{+} \nt{T} \nt{E'} \\
\texttt{+} \nt{T} \nt{E'} \texttt{)} \nt{T'} \nt{E'} \texttt{;} \nt{lista\_sent} \texttt{\$} & \texttt{+} \texttt{num} \texttt{;} \texttt{\$} & pop/avanza \\
\nt{T} \nt{E'} \texttt{)} \nt{T'} \nt{E'} \texttt{;} \nt{lista\_sent} \texttt{\$} & \texttt{num} \texttt{;} \texttt{\$} & \nt{T} $\rightarrow$ \nt{F} \nt{T'} \\
\nt{F} \nt{T'} \nt{E'} \texttt{)} \nt{T'} \nt{E'} \texttt{;} \nt{lista\_sent} \texttt{\$} & \texttt{num} \texttt{;} \texttt{\$} & \nt{F} $\rightarrow$ \texttt{num} \\
\texttt{num} \nt{T'} \nt{E'} \texttt{)} \nt{T'} \nt{E'} \texttt{;} \nt{lista\_sent} \texttt{\$} & \texttt{num} \texttt{;} \texttt{\$} & pop/avanza \\
\nt{T'} \nt{E'} \texttt{)} \nt{T'} \nt{E'} \texttt{;} \nt{lista\_sent} \texttt{\$} & \texttt{;} \texttt{\$} & \nt{T'} $\rightarrow$ $\varepsilon$ \\
\nt{E'} \texttt{)} \nt{T'} \nt{E'} \texttt{;} \nt{lista\_sent} \texttt{\$} & \texttt{;} \texttt{\$} & \nt{E'} $\rightarrow$ $\varepsilon$ \\
\texttt{)} \nt{T'} \nt{E'} \texttt{;} \nt{lista\_sent} \texttt{\$} & \texttt{;} \texttt{\$} & Error: se esperaba \texttt{)} \\
\bottomrule
\end{tabular}
\caption{Simulación de \cmd{x = (10 + 5;}.}
\label{tab:sim2}
\end{table}

En la última fila el tope es el terminal \cmd{)} y la cadena empieza con \cmd{;}: no hay
\emph{pop/avanza} posible. Es exactamente el error que reporta el programa:
\cmd{Se esperaba ')' antes de ';'} en la columna 12 (figura~\ref{fig:captura_06_parentesis.png}).

% ===========================================================================
\section{Implementación}
\label{sec:implementacion}

\needspace{8\baselineskip}
\subsection{Arquitectura}

El programa tiene cuatro archivos fuente. \cmd{main.c} abre el archivo, llama al analizador
léxico, que deja todos los tokens en un arreglo, y después al sintáctico, que solo recorre ese
arreglo. Ninguno de los dos analizadores imprime el resultado final: eso lo hace \cmd{main.c}
con los conteos de error que cada uno devuelve.

\begin{figure}[!htbp]
\centering
\resizebox{\linewidth}{!}{%
\begin{tikzpicture}[font=\small, node distance=0.55cm and 0.6cm, >=Stealth,
  caja/.style={draw, rectangle, minimum height=1.0cm, align=center, inner sep=5pt},
  dato/.style={draw, rectangle, rounded corners=6pt, minimum height=0.9cm, align=center,
               fill=grisclaro, inner sep=5pt}]
  \node[dato] (fuente) {Archivo\\\cmd{.txt}};
  \node[caja, right=of fuente] (lex) {\cmd{analisis\_lexico()}\\\cmd{lexico.l} (Flex)};
  \node[dato, right=of lex] (buf) {\cmd{Token tokens[]}\\termina en \cmd{FIN}};
  \node[caja, right=of buf] (sin) {\cmd{analisis\_sintactico()}\\\cmd{sintactico.c}};
  \node[dato, below=of lex] (ts) {Tabla de\\símbolos};
  \node[caja, below=of sin] (res) {Resultado en\\\cmd{main.c}};
  \draw[->] (fuente) -- (lex);
  \draw[->] (lex) -- (buf);
  \draw[->] (buf) -- (sin);
  \draw[->] (lex) -- (ts);
  \draw[->] (sin) -- node[right, font=\footnotesize] {errores} (res);
  \draw[->] (lex.south east) |- node[pos=0.75, below, font=\footnotesize] {errores léxicos} (res.west);
\end{tikzpicture}}
\caption{Flujo de datos entre los módulos.}
\end{figure}

Los listados de esta sección son el código del repositorio, con los mismos números de línea que
los archivos.

\needspace{8\baselineskip}
\subsection{Interfaz común: \texttt{tokens.h}}

\lstinputlisting[style=fuente, caption={\texttt{src/tokens.h}}, label=cod:tokens]{../src/tokens.h}

\begin{itemize}[leftmargin=1.1cm,itemsep=0.25em]
\item \textbf{Líneas 7 a 14.} \cmd{TipoToken} enumera las clases léxicas con los nombres de la
asignación. \cmd{FIN} vale 0 porque \cmd{yylex()} devuelve 0 al llegar al fin de archivo, así
que el fin de entrada se convierte solo en el símbolo \cmd{\$}.
\item \textbf{Líneas 16 a 21.} Cada \cmd{Token} guarda su clase, una copia del lexema y la línea
y columna donde empieza. Con eso los dos analizadores pueden reportar la posición de cualquier
error.
\item \textbf{Líneas 23 a 30.} Son las únicas funciones que un módulo ve del otro.
\cmd{analisis\_lexico()} devuelve el número de tokens; \cmd{analisis\_sintactico()}, el número
de errores sintácticos.
\end{itemize}

\needspace{8\baselineskip}
\subsection{Analizador léxico: \texttt{lexico.l}}
\label{sec:lexico}

\needspace{6\baselineskip}
\subsubsection*{Sección de definiciones}

\lstinputlisting[style=fuente, firstline=1, lastline=40, firstnumber=1,
caption={\texttt{src/lexico.l}, definiciones}, label=cod:lex-def]{../src/lexico.l}

\begin{itemize}[leftmargin=1.1cm,itemsep=0.25em]
\item \textbf{Líneas 11 a 15.} Variables de estado del scanner: la posición del siguiente
carácter (\cmd{linea}, \cmd{columna}), la posición donde empezó el lexema actual
(\cmd{tok\_linea}, \cmd{tok\_columna}), dónde se abrió un comentario de bloque, el contador de
errores y la bandera \cmd{sustituido}, que marca un lexema con error entregado como token válido.
\item \textbf{Líneas 21 a 26.} \cmd{YY\_USER\_ACTION} se ejecuta antes de la acción de cada
regla. Primero guarda la posición de inicio y luego recorre \cmd{yytext}: un salto de línea
incrementa \cmd{linea} y regresa la columna a 1; cualquier otro byte avanza la columna, excepto
los bytes de continuación UTF-8 (\cmd{10xxxxxx}), para que \cmd{ñ} cuente como un solo carácter.
\item \textbf{Línea 29.} \cmd{noyywrap} indica que solo hay un archivo de entrada;
\cmd{nounput} y \cmd{noinput} quitan funciones que no se usan y evitan advertencias de
\cmd{gcc}; \cmd{8bit} permite leer bytes mayores que 127.
\item \textbf{Línea 30.} \cmd{COMENTARIO} es un estado exclusivo: mientras está activo solo se
prueban las reglas que empiezan con \cmd{<COMENTARIO>}.
\item \textbf{Líneas 32 a 40.} Las definiciones regulares de la sección~\ref{sec:tokens}.
\end{itemize}

\needspace{6\baselineskip}
\subsubsection*{Sección de reglas}

\lstinputlisting[style=fuente, firstline=42, lastline=88, firstnumber=42,
caption={\texttt{src/lexico.l}, reglas}, label=cod:lex-reglas]{../src/lexico.l}

\begin{itemize}[leftmargin=1.1cm,itemsep=0.25em]
\item \textbf{Líneas 44 y 45.} Los espacios y el comentario de línea tienen acción vacía: se
reconocen y se descartan, pero \cmd{YY\_USER\_ACTION} ya actualizó línea y columna.
\item \textbf{Líneas 46 a 53.} \cmd{/*} guarda la posición de apertura y entra a
\cmd{COMENTARIO}. Ahí \cmd{*/} regresa a \cmd{INITIAL} y cualquier otro carácter se ignora. Si
se llega al fin de archivo dentro del comentario, se reporta el error en la línea y columna de
apertura y se devuelve \cmd{FIN}.
\item \textbf{Líneas 55 a 62.} Las palabras reservadas van antes de \cmd{\{ID\}} para ganar los
empates de longitud. Cada acción solo devuelve la clase; el lexema queda en \cmd{yytext}.
\item \textbf{Líneas 66 a 68.} Los lexemas mal formados se reportan y se devuelven como
\cmd{NUMBER} o \cmd{ID}, marcando \cmd{sustituido} para que no entren a la tabla de símbolos.
\item \textbf{Líneas 70 a 86.} Operadores y delimitadores; los de dos caracteres van primero.
\item \textbf{Línea 88.} Última regla: un carácter UTF-8 o cualquier byte que no reconoció
ninguna regla anterior es un carácter no válido. No devuelve token, así que el carácter se
descarta y el análisis sigue.
\end{itemize}

\needspace{6\baselineskip}
\subsubsection*{Código de usuario}

\lstinputlisting[style=fuente, firstline=90, lastline=163, firstnumber=90,
caption={\texttt{src/lexico.l}, código de usuario}, label=cod:lex-codigo]{../src/lexico.l}

\begin{itemize}[leftmargin=1.1cm,itemsep=0.25em]
\item \textbf{Líneas 92 a 99.} \cmd{nombres[]} tiene el texto de cada clase en el mismo orden que
el \cmd{enum}; \cmd{nombre\_token()} lo usa para imprimir la tabla de tokens.
\item \textbf{Líneas 103 a 110.} \cmd{error\_lexico()} cuenta el error e imprime el mensaje con
el formato \cmd{Error léxico en línea L, columna C: descripción 'lexema'.}
\item \textbf{Líneas 112 a 118.} \cmd{crecer()} duplica la capacidad de un arreglo dinámico; lo
usan el buffer de tokens y la tabla de símbolos.
\item \textbf{Líneas 120 a 139.} Tabla de símbolos: un arreglo de \cmd{Simbolo} con nombre,
línea de la primera aparición y número de apariciones. \cmd{registrar\_simbolo()} busca el nombre
de forma lineal; si existe, suma una aparición, y si no, lo agrega.
\item \textbf{Líneas 142 a 163.} \cmd{analisis\_lexico()} es la interfaz con el resto del
programa. Llama a \cmd{yylex()} hasta recibir \cmd{FIN}, guarda cada token con su lexema y
posición, y registra los identificadores válidos. El último elemento es \cmd{FIN} con lexema
\cmd{\$}. Como el arreglo está completo antes del análisis sintáctico, \cmd{main.c} puede
imprimir la tabla de tokens primero, igual que el arreglo \cmd{tokens[]} del ejemplo en C del
tema 4.4.
\end{itemize}

\needspace{8\baselineskip}
\subsection{Analizador sintáctico: \texttt{sintactico.c}}
\label{sec:sintactico}

\needspace{6\baselineskip}
\subsubsection*{Estado, \texttt{match()} y reporte de errores}

\lstinputlisting[style=fuente, firstline=10, lastline=71, firstnumber=10,
caption={\texttt{src/sintactico.c}, estado y funciones de apoyo}, label=cod:sin-apoyo]{../src/sintactico.c}

\begin{itemize}[leftmargin=1.1cm,itemsep=0.25em]
\item \textbf{Líneas 10 a 14.} \cmd{lookahead} apunta al token actual del buffer; \cmd{i} es su
índice. \cmd{panico} vale 1 desde que se detecta un error hasta que se sincroniza.
\item \textbf{Líneas 17 a 24.} \cmd{esperado[]} da el texto de cada terminal para los mensajes,
por ejemplo \cmd{')'} o \cmd{identificador}.
\item \textbf{Líneas 26 a 45.} Prototipos: una función por no terminal de la gramática.
\item \textbf{Líneas 47 a 50.} \cmd{avanzar()} mueve el lookahead al siguiente token, sin pasar
de \cmd{FIN}.
\item \textbf{Líneas 55 a 64.} \cmd{error()} imprime la línea y columna del token donde se
detectó el error, lo que se esperaba y lo que se encontró. Si ya está en pánico no imprime nada,
así un solo error no produce varios mensajes.
\item \textbf{Líneas 66 a 71.} \cmd{match(t)} es el \emph{pop/avanza} de la simulación: si el
lookahead es de la clase \cmd{t}, avanza; si no, reporta que se esperaba \cmd{t}.
\end{itemize}

\needspace{6\baselineskip}
\subsubsection*{Recuperación en modo pánico}

\lstinputlisting[style=fuente, firstline=73, lastline=93, firstnumber=73,
caption={\texttt{src/sintactico.c}, sincronización}, label=cod:sin-panico]{../src/sintactico.c}

\cmd{sincronizar()} descarta tokens hasta el conjunto \{\cmd{;}, \cmd{\}}, \cmd{\$}\}. La
variable \cmd{nivel} cuenta las llaves abiertas en el camino: si un error ocurre en el encabezado
de un \cmd{while}, su bloque se salta completo y no queda un \cmd{\}} suelto (líneas 86 a 88). Un
\cmd{\}} con \cmd{nivel} en cero no se consume (línea 84) porque cierra el bloque que contiene al
error. Con \cmd{era\_if}, el \cmd{else} de un \cmd{if} con error también se salta.

\needspace{6\baselineskip}
\subsubsection*{Programa y sentencias}

\lstinputlisting[style=fuente, firstline=95, lastline=127, firstnumber=95,
caption={\texttt{src/sintactico.c}, programa y sentencias}, label=cod:sin-sent]{../src/sintactico.c}

\begin{itemize}[leftmargin=1.1cm,itemsep=0.25em]
\item \textbf{Líneas 96 a 105.} \cmd{programa()} analiza la lista de sentencias. Si al regresar
el lookahead no es \cmd{FIN}, encontró un \cmd{\}} sin \cmd{\{}: lo reporta, lo descarta y sigue.
\item \textbf{Líneas 108 a 115.} \cmd{lista\_sentencias()} aplica la producción 3
($\varepsilon$) cuando el lookahead es \cmd{\}} o \cmd{\$}, que es Follow(\nt{lista\_sent}). En
otro caso aplica la producción 2. Aquí se llama a \cmd{sincronizar()} cuando la sentencia
terminó en pánico, de modo que el error queda acotado a una sentencia.
\item \textbf{Líneas 118 a 127.} \cmd{sentencia()} es la fila \nt{sentencia} de la tabla
LL(1): el \cmd{switch} elige la producción 4, 5, 6 o 7 según el lookahead, y en las celdas
$\nexists$ reporta el error.
\end{itemize}

\needspace{6\baselineskip}
\subsubsection*{Declaraciones y asignación}

\lstinputlisting[style=fuente, firstline=129, lastline=176, firstnumber=129,
caption={\texttt{src/sintactico.c}, declaraciones y asignación}, label=cod:sin-decl]{../src/sintactico.c}

Cada función copia el lado derecho de su producción: un terminal es un \cmd{match()} y un no
terminal, una llamada. \cmd{lista\_id()} e \cmd{inicializacion()} tienen producción
$\varepsilon$: solo entran a la alternativa larga si el lookahead es \cmd{,} o \cmd{=}; en otro
caso no consumen nada y el \cmd{;} lo verifica \cmd{declaracion()}.

\needspace{6\baselineskip}
\subsubsection*{Estructuras de control}

\lstinputlisting[style=fuente, firstline=178, lastline=222, firstnumber=178,
caption={\texttt{src/sintactico.c}, if, while y bloques}, label=cod:sin-control]{../src/sintactico.c}

\cmd{parte\_else()} decide entre las producciones 17 y 18 con un solo token: si sigue \cmd{else}
lo consume, y si no, aplica $\varepsilon$. \cmd{resto\_else()} distingue \cmd{else if} de
\cmd{else \{} por el token siguiente (producciones 19 y 20). \cmd{bloque()} reutiliza
\cmd{lista\_sentencias()}, así que los bloques pueden anidarse a cualquier profundidad.

\needspace{6\baselineskip}
\subsubsection*{Condiciones y expresiones}

\lstinputlisting[style=fuente, firstline=224, lastline=293, firstnumber=224,
caption={\texttt{src/sintactico.c}, condiciones y expresiones}, label=cod:sin-expr]{../src/sintactico.c}

Son las funciones \cmd{E}, \cmd{E\_Prime}, \cmd{T}, \cmd{T\_Prime} y \cmd{F} del tema 4.4.
\cmd{E\_Prime()} y \cmd{T\_Prime()} terminan sin hacer nada cuando el lookahead no es un
operador de su nivel: es la producción $\varepsilon$, y el token queda para la función que las
llamó. \cmd{F()} tiene un caso por cada producción 38 a 41 y un \cmd{default} que corresponde a
las celdas $\nexists$ de su fila.

\needspace{6\baselineskip}
\subsubsection*{Punto de entrada}

\lstinputlisting[style=fuente, firstline=295, lastline=303, firstnumber=295,
caption={\texttt{src/sintactico.c}, punto de entrada}, label=cod:sin-entrada]{../src/sintactico.c}

Recibe el buffer, coloca el lookahead en el primer token, analiza \nt{programa} y devuelve
cuántos errores sintácticos encontró.

\needspace{8\baselineskip}
\subsection{Programa principal: \texttt{main.c}}

\lstinputlisting[style=fuente, caption={\texttt{src/main.c}}, label=cod:main]{../src/main.c}

\begin{itemize}[leftmargin=1.1cm,itemsep=0.25em]
\item \textbf{Líneas 6 a 14.} Valida que se reciba un archivo y que se pueda abrir; en otro caso
termina con código 2.
\item \textbf{Líneas 16 a 19.} Ejecuta el análisis léxico completo. Los errores léxicos se
imprimen en este momento, por eso aparecen antes de la tabla de tokens.
\item \textbf{Líneas 21 a 25.} Imprime la tabla de tokens con línea, columna, clase y lexema.
\item \textbf{Líneas 27 a 36.} Imprime la misma secuencia agrupada por línea, en el formato del
enunciado (\cmd{INT ID SEMICOLON}). \cmd{FIN} no se muestra.
\item \textbf{Líneas 38 a 43.} Imprime la tabla de símbolos y ejecuta el análisis sintáctico.
\item \textbf{Líneas 45 a 52.} Indica si el programa es válido y termina con código 0 (válido) o
1 (con errores), lo que permite usarlo desde el script de pruebas.
\end{itemize}

% ===========================================================================
\section{Manejo de errores}
\label{sec:errores}

\begin{respuesta}
\small\ttfamily
Error léxico en línea L, columna C: descripción 'lexema'.\\
Error sintáctico en línea L, columna C: Se esperaba X antes de 'Y'.
\end{respuesta}

\begin{table}[!htbp]
\centering\small
\begin{tabularx}{\textwidth}{@{}>{\raggedright\arraybackslash}p{4.2cm} l X@{}}
\toprule
\textbf{Error léxico} & \textbf{Ejemplo} & \textbf{Recuperación} \\
\midrule
Carácter no válido & \cmd{@ \# \$ ñ} & Se descarta y el análisis sigue \\
Número mal formado & \cmd{3.14.5}, \cmd{7.} & Se entrega como \cmd{NUMBER} \\
Identificador mal formado & \cmd{2abc}, \cmd{año} & Se entrega como \cmd{ID} \\
Comentario sin cerrar & \cmd{/* \dots} & Se reporta en la apertura; termina la entrada \\
\bottomrule
\end{tabularx}
\caption{Errores léxicos y su recuperación.}
\end{table}

Entregar los lexemas mal formados con la clase que se intentaba escribir evita errores
sintácticos en cascada: en \cmd{x = 3.14.5;} se reporta el número y la asignación sigue siendo
correcta.

Un error sintáctico se detecta cuando \cmd{match()} no encuentra el terminal esperado o cuando
una función cae en una celda $\nexists$ de su fila. Cuando hay varias alternativas válidas se
enumeran (\cmd{Se esperaba identificador, número, '(' o '-' antes de ')'}). La recuperación es
en modo pánico (temas 4.3 y 4.4):

\begin{enumerate}[leftmargin=1.1cm,itemsep=0.25em]
\item Se activa la bandera \cmd{panico}. Mientras está activa no se reportan más errores,
\cmd{match()} no consume tokens y las funciones regresan hasta \cmd{lista\_sentencias()}.
\item \cmd{sincronizar()} descarta tokens hasta \{\cmd{;}, \cmd{\}}, \cmd{\$}\}, saltando los
bloques completos (código~\ref{cod:sin-panico}).
\item Se desactiva la bandera y el análisis sigue con la siguiente sentencia.
\end{enumerate}

% ===========================================================================
\section{Funcionamiento con un caso de ejemplo}
\label{sec:ejemplo}

Esta sección sigue un solo caso, el programa de ejemplo de la asignación, desde la compilación
hasta el resultado. Las figuras son capturas de una terminal de Linux con los comandos
ejecutados en la carpeta del repositorio.

\needspace{8\baselineskip}
\subsection{Compilación}

\captura{captura_01_compilacion.png}{Compilación del proyecto con \cmd{make}.}

\cmd{make} ejecuta tres pasos: crea la carpeta \cmd{build}, genera \cmd{build/lex.yy.c} a partir
de \cmd{src/lexico.l} con Flex y compila ese archivo junto con \cmd{sintactico.c} y
\cmd{main.c}. \cmd{gcc} se ejecuta con \cmd{-Wall -Wextra} y no muestra ninguna advertencia.

\needspace{8\baselineskip}
\subsection{Programa de entrada}

\captura{captura_02_entrada.png}{Programa de ejemplo de la asignación
(\cmd{pruebas/validos/01\_ejemplo\_enunciado.txt}).}

El programa declara dos variables, hace dos asignaciones, un \cmd{if} y un \cmd{while}, así que
usa todas las construcciones mínimas que pide la asignación.

\needspace{8\baselineskip}
\subsection{Ejecución y lectura de la salida}

\captura[height=0.84\textheight]{captura_03_ejecucion.png}{Salida completa del analizador para el
programa de ejemplo.}

La salida de la figura~\ref{fig:captura_03_ejecucion.png} se lee de arriba hacia abajo y cada
bloque sale de una parte del código:

\begin{enumerate}[leftmargin=1.1cm,itemsep=0.3em]
\item \textbf{Análisis léxico.} No hay mensajes de error, porque ninguna regla de error de
\cmd{lexico.l} (líneas 66 a 68 y 88) se activó.
\item \textbf{Tokens reconocidos.} Una fila por token, impresa por \cmd{main.c} (líneas 21 a
25). La columna sale de \cmd{YY\_USER\_ACTION}: en la línea 1, \cmd{int} empieza en la columna
1, \cmd{x} en la 5 y \cmd{;} en la 6.
\item \textbf{Secuencia de tokens por línea.} La misma información agrupada por línea
(\cmd{main.c}, líneas 27 a 36). Coincide con la secuencia esperada en el enunciado:
\cmd{INT ID SEMICOLON}, \cmd{ID ASSIGN ID PLUS NUMBER SEMICOLON},
\cmd{IF LPAREN ID GREATER NUMBER RPAREN LBRACE}, etcétera.
\item \textbf{Tabla de símbolos.} \cmd{x} aparece por primera vez en la línea 1 y se usa 7 veces;
\cmd{y}, en la línea 2 y 4 veces. Los cuenta \cmd{registrar\_simbolo()} (\cmd{lexico.l}, líneas
125 a 131).
\item \textbf{Análisis sintáctico.} Las funciones recorrieron todo el buffer sin llamar a
\cmd{error()}. Para la línea 4, \cmd{y = x + 5;}, el recorrido es el de la
tabla~\ref{tab:sim1} y el árbol de la figura~\ref{fig:arbol}.
\item \textbf{Resultado.} Sin errores de ninguna fase, el programa es válido y termina con
código 0.
\end{enumerate}

\needspace{8\baselineskip}
\subsection{El mismo programa con errores}

Para ver el manejo de errores se modificó el ejemplo con tres errores: un carácter no válido en
la línea 4, un paréntesis faltante en la línea 5 y un punto y coma faltante en la línea 9.

\captura{captura_04_errores_entrada.png}{Programa de ejemplo con tres errores
(\cmd{docs/evidencias/ejemplo\_con\_errores.txt}).}

\captura[height=0.84\textheight]{captura_05_errores_ejecucion.png}{Salida del analizador para el
programa con errores.}

\begin{enumerate}[leftmargin=1.1cm,itemsep=0.3em]
\item \textbf{Línea 4, columna 11: \cmd{@}.} Lo detecta la regla de la línea 88 de
\cmd{lexico.l} durante el análisis léxico, por eso el mensaje aparece antes de la tabla de
tokens. El carácter se descarta y la línea 4 queda como \cmd{ID ASSIGN ID PLUS NUMBER SEMICOLON},
una asignación correcta: no hay error sintáctico en cascada.
\item \textbf{Línea 5, columna 12: falta \cmd{)}.} En la secuencia por línea se ve
\cmd{NUMBER LBRACE} donde debería ir \cmd{NUMBER RPAREN LBRACE}. \cmd{sent\_if()} llama a
\cmd{match(RPAREN)} con el lookahead en \cmd{\{} y se reporta
\cmd{Se esperaba ')' antes de '\{'}. En modo pánico, \cmd{sincronizar()} salta el bloque del
\cmd{if} completo (líneas 5 a 7).
\item \textbf{Línea 10, columna 1: falta \cmd{;}.} La asignación de la línea 9 termina en
\cmd{NUMBER} y el siguiente token es el \cmd{\}} de la línea 10, que es donde se detecta el
error. El analizador sincroniza en ese \cmd{\}} y cierra el \cmd{while} normalmente.
\item \textbf{Resultado.} Un error léxico y dos sintácticos; el programa no es válido y termina
con código 1.
\end{enumerate}

\needspace{8\baselineskip}
\subsection{Caso de error de la asignación}

\captura{captura_06_parentesis.png}{Entrada \cmd{x = (10 + 5;} y su salida.}

El mensaje coincide con el de la asignación: \cmd{Error sintáctico en línea 1, columna 12: Se
esperaba ')' antes de ';'.} La columna 12 es la del \cmd{;}, el primer token que no puede seguir
a la expresión dentro del paréntesis. La tabla~\ref{tab:sim2} termina en el mismo punto.

\needspace{8\baselineskip}
\subsection{Batería de pruebas}

\cmd{pruebas/validos} tiene seis programas que deben aceptarse y \cmd{pruebas/invalidos} ocho que
deben rechazarse. Cada inválido indica en un comentario inicial los errores que contiene.
\cmd{make pruebas} los ejecuta todos, muestra una línea por caso y guarda la salida completa de
cada uno en \cmd{pruebas/resultados/}.

\captura{captura_07_pruebas.png}{Ejecución de la batería de pruebas.}

\begin{table}[!htbp]
\centering\small
\begin{tabular}{@{}l l c c@{}}
\toprule
\textbf{Archivo} & \textbf{Resultado} & \textbf{Léxicos} & \textbf{Sintácticos} \\
\midrule
\cmd{validos/01\_ejemplo\_enunciado} & Aceptado & 0 & 0 \\
\cmd{validos/02\_declaraciones\_y\_comentarios} & Aceptado & 0 & 0 \\
\cmd{validos/03\_expresiones\_aritmeticas} & Aceptado & 0 & 0 \\
\cmd{validos/04\_condicionales} & Aceptado & 0 & 0 \\
\cmd{validos/05\_ciclos\_anidados} & Aceptado & 0 & 0 \\
\cmd{validos/06\_programa\_completo} & Aceptado & 0 & 0 \\
\midrule
\cmd{invalidos/01\_lexico\_caracteres\_invalidos} & Rechazado & 4 & 2 \\
\cmd{invalidos/02\_lexico\_literales\_mal\_formados} & Rechazado & 3 & 0 \\
\cmd{invalidos/03\_lexico\_comentario\_sin\_cerrar} & Rechazado & 1 & 0 \\
\cmd{invalidos/04\_sintactico\_parentesis} & Rechazado & 0 & 3 \\
\cmd{invalidos/05\_sintactico\_punto\_y\_coma} & Rechazado & 0 & 3 \\
\cmd{invalidos/06\_sintactico\_estructuras\_control} & Rechazado & 0 & 4 \\
\cmd{invalidos/07\_sintactico\_declaraciones} & Rechazado & 0 & 6 \\
\cmd{invalidos/08\_mixto\_lexico\_y\_sintactico} & Rechazado & 2 & 3 \\
\bottomrule
\end{tabular}
\caption{Resultado de la batería de pruebas.}
\end{table}

% ===========================================================================
\needspace{12\baselineskip}
\section{Instalación y ejecución}

Requiere \cmd{gcc}, \cmd{flex} y \cmd{make}, sin otras dependencias. En Linux o WSL:

\begin{lstlisting}
sudo apt install build-essential flex # instala las herramientas
make                                  # compila
./analizador programa.txt             # analiza un archivo
make pruebas                          # corre los 14 casos de prueba
\end{lstlisting}

El programa termina con código 0 si la entrada es válida, 1 si tiene errores y 2 si no se pudo
abrir el archivo.

\begin{table}[!htbp]
\centering\small
\begin{tabularx}{\textwidth}{@{}l X@{}}
\toprule
\textbf{Ruta} & \textbf{Contenido} \\
\midrule
\cmd{src/} & Código fuente de los dos analizadores y del programa principal \\
\cmd{pruebas/} & Casos válidos e inválidos, script de ejecución y salidas \\
\cmd{docs/} & Este documento, su fuente LaTeX y las capturas de la sección~\ref{sec:ejemplo} \\
\cmd{Makefile} & Reglas \cmd{make}, \cmd{make pruebas}, \cmd{make docs} y \cmd{make clean} \\
\bottomrule
\end{tabularx}
\caption{Organización del repositorio.}
\end{table}

% ===========================================================================
\section{Limitaciones}

\begin{itemize}[leftmargin=1.1cm,itemsep=0.25em]
\item No hay análisis semántico: no se verifica la declaración previa de variables ni la
compatibilidad de tipos.
\item Una condición admite una sola comparación; el lenguaje no tiene operadores lógicos.
\item El modo pánico puede descartar una sentencia correcta que sigue a un error: si falta el
\cmd{;} de una línea, la sincronización llega al \cmd{;} de la siguiente.
\end{itemize}

% ===========================================================================
\section*{Referencias}
\addcontentsline{toc}{section}{Referencias}

\begin{itemize}[leftmargin=1.2cm,itemsep=0.5em]
\item Aho, A. V., Lam, M. S., Sethi, R. y Ullman, J. D. \emph{Compiladores: principios, técnicas
y herramientas}, 2.ª ed. Pearson, 2008.
\item Castañeda Castañeda, M. E. \emph{Presentaciones de la asignatura Compiladores}, temas 1.4,
2.1 a 2.6, 3.1, 3.2 y 4.1 a 4.4. Facultad de Ingeniería, UNAM, 2026.
\item Paxson, V., Estes, W. y Millaway, J. \emph{Lexical Analysis With Flex, for Flex 2.6.4}.
Manual de Flex.
\end{itemize}

% ===========================================================================
\appendix
\renewcommand{\thesection}{\Alph{section}}
\section{Declaración de uso de inteligencia artificial}

La asignatura permite usar herramientas de inteligencia artificial si se declaran las
instrucciones empleadas. En este proyecto se usó un asistente de IA generativa con acceso al
navegador, que consultó en Classroom la asignación, la rúbrica y las presentaciones del curso. El
asistente generó el código fuente, los casos de prueba y este documento, y preparó el historial
de versiones del repositorio. El equipo revisó el contenido, ejecutó las pruebas y se encarga de
la demostración.

Las instrucciones que definieron el contenido técnico del proyecto se presentan en el orden en
que se dieron. La redacción está resumida y editada, y conserva el sentido de cada una.

\begin{apartado}{INSTRUCCIÓN 1 \ \textbullet\ \ DESARROLLO DEL PROYECTO}
Desarrolla el proyecto final de Compiladores conforme a la asignación y la rúbrica publicadas en
Classroom. Toma de las presentaciones del profesor los métodos vistos en clase (expresiones
regulares con Flex, gramáticas LL(1), análisis descendente recursivo y recuperación en modo
pánico) y aplícalos sin sustituirlos por otras técnicas. Escribe código sencillo y fácil de
mantener, y documentación con un registro técnico profesional. Publica el proyecto en un
repositorio de GitHub con un commit por etapa e invita a los demás integrantes.
\end{apartado}

\begin{apartado}{INSTRUCCIÓN 2 \ \textbullet\ \ RECOPILACIÓN Y PLANEACIÓN}
Antes de programar, reúne la asignación, la rúbrica y el material del curso, y redacta documentos
de requisitos y de planeación que sirvan de contexto para el resto del trabajo. El repositorio se
crea al final.
\end{apartado}

\begin{apartado}{INSTRUCCIÓN 3 \ \textbullet\ \ ALCANCE DE LA ENTREGA}
Implementa el proyecto completo: analizadores, pruebas y documento técnico. Guarda el contexto
del trabajo en archivos Markdown dentro de un archivo comprimido. Los integrantes son Andrés
Basilio Illescas, Samuel Santiago Cruz Macedo y Jesús Jafet Fernández Hernandez.
\end{apartado}

\begin{apartado}{INSTRUCCIÓN 4 \ \textbullet\ \ REVISIÓN CONTRA LA RÚBRICA}
Revisa el estado del proyecto frente a la rúbrica, identifica lo que falta y de qué depende, y
haz los ajustes necesarios para obtener la calificación máxima.
\end{apartado}

\begin{apartado}{INSTRUCCIÓN 5 \ \textbullet\ \ AJUSTES DE FORMATO Y EJECUCIÓN}
Adapta el documento técnico al formato LaTeX que el equipo usa en sus reportes y reduce su
extensión. Simplifica la compilación y la ejecución de las pruebas para que el proyecto se corra
con pocos comandos. Crea el repositorio en GitHub y envía las invitaciones a los integrantes.
\end{apartado}

\begin{apartado}{INSTRUCCIÓN 6 \ \textbullet\ \ CONGRUENCIA CON EL CURSO Y EVIDENCIAS}
Verifica que el proyecto use los métodos y la notación de las tareas de clase (First y Follow,
tabla LL(1) con celdas sin correspondencia, simulación con pila). Usa los logos institucionales
en blanco y negro. Documenta el código parte por parte y muestra, con capturas de la ejecución, un
caso de ejemplo completo con su salida y sus errores.
\end{apartado}

\end{document}
